\documentclass{article}

\usepackage[utf8]{inputenc}
\usepackage[T1]{fontenc}
\usepackage{lmodern}
\usepackage{amsmath}
\usepackage{amssymb}
\usepackage{amsthm}
\usepackage{mathtools}
\usepackage{empheq}
\usepackage{parskip}
\usepackage{cleveref}
\usepackage[margin=0.75in]{geometry}

\title{MATH1013 Assignment 2}
\author{London Wallace | 23071503}
\date{}

\setlength{\parindent}{0pt}
\setlength{\parskip}{1em plus 0.1em minus 0.1em}

\begin{document}

\maketitle{}

\begin{flushleft}
	\textbf{Due Date:} April 14, 2026, 11:59 PM
\end{flushleft}

\section*{Theorems}
Herein, \textit{Calculus} and \textit{Spivak} will refer to the book \textit{Calculus, Third Edition} by
\textit{Michael Spivak} published in 1994, printed in 2009. \textbf{ISBN 978-0-521-86744-3}.

All theorems published in \textit{Spivak} and relied upon for this assignment are printed below in order of appearance.

\begin{enumerate}
	\item (Theorem 10-4, pg. 167) If $f$ and $g$ are differentiable at $a$, then
	      \begin{equation}
		      (f \cdot g)'(a)=f'(a)g(a)+g'(a)f(a) \label{product-rule}
	      \end{equation}
	\item (Theorem 10-9, pg. 176) If $g$ is differentiable at $a$, and $f$ is differentiable at $g(a)$, then
	      $f \circ g$ is differentiable at $a$, and
	      \begin{equation}
		      (f \circ g)'(a)=f'(g(a))\cdot g'(a) \label{chain-rule}
	      \end{equation}
	\item (Theorem 10-5, pg. 167) If $g(x)=c \cdot f(x)$, and $f$ is differentiable at $a$, then $g$ is differentiable
	      at $a$, and
	      \begin{equation}
		      g'(a)=c \cdot f'(a) \label{product-rule-for-constants}
	      \end{equation}
	\item (Theorem 10-7, pg. 169) If $g$ is differentiable at $a$, and $g(a) \ne 0$, then $\frac1g$ is differentiable at $a$, and
	      \begin{equation}
		      \left( \dfrac{1}{g} \right)'(a)=\dfrac{-g'(a)}{(g(a))^2} \label{ddx-of-1on-g}
	      \end{equation}
	\item (Theorem 5-2, pg. 102) If $\lim_{x \to a}f(x)=l$ and $\lim_{x \to a}g(x)=m$, then
	      \begin{equation}
		      \lim_{x \to a}(f+g)(x) = l+ m \label{limit-of-sums}
	      \end{equation}
	      \begin{equation}
		      \lim_{x \to x}(f \cdot g)(x) =l \cdot m	 \label{limit-of-products}
	      \end{equation}
	\item (Theorem 6-1, pg. 114) If $f$ and $g$ are continuous at $a$, then
	      \begin{equation}
		      f+g \text{ is continuous at }a \label{continuous-sum-is-continuous}
	      \end{equation}
	\item (Theorem 7-1, pg. 120) If $f$ is continuous on $[a,b]$ and $f(a) < 0 < f(b)$, then there is some $x \in [a,b]$ such that
	      \begin{equation}
		      f(x)=0 \label{mean-value}
	      \end{equation}
	\item (Theorem 11-9, pg. 201) Suppose that $\lim_{x \to a}f(x)=0$, and $\lim_{x \to a}g(x)=0$, and suppose that the $\lim_{x \to a}\dfrac{f'(x)}{g'(x)}$
	      exists. Then the $\lim_{x \to a}\dfrac{f(x)}{g(x)}$ exists and
	      \begin{equation}
		      \lim_{x \to a}\dfrac{f(x)}{g(x)}=\lim_{x \to a}\dfrac{f'(x)}{g'(x)} \label{le-medical-centre}
	      \end{equation}
	\item (Theorem 10-6, pg. 168) If $f(x)=x^n$ for some natural number $n$, then
	      \begin{equation}
		      f'(a)=na^{n-1} \quad \forall a \label{power-rule}
	      \end{equation}
\end{enumerate}


\section{Exercise 1}
It is given that $\forall x \in \mathbb{R}$ the first derivative $(\cos(x))' = -\sin(x)$.

For $\alpha, \beta \in (0,+\infty)$ let $f : [0,1) \to \mathbb{R}$ be defined:

\begin{equation}
	f(x) \coloneqq
	\begin{cases}
		x^{\alpha}\cos \left(\dfrac{1}{x^\beta}\right) & x \in (0,1) \\
		0                                              & x = 0
	\end{cases}
	\label{exercise1}
\end{equation}

\subsection{}
Compute $f'(x)$ for every $x \in (0,1)$

\subsubsection{Solution}

For $x \in (0,1)$, $f(x)=x^{\alpha}\cos \left(\dfrac{1}{x^\beta}\right)$ from \cref{exercise1}.

Define $g(x)\coloneq x^{\alpha}$. Define $h(x)\coloneq\left(\dfrac{1}{x^{\beta}}\right)$.

Then,
\begin{equation}
	f(x) = g(x) \cos(h(x)) \label{exercise1.1.1}
\end{equation}
Thus, from \cref{product-rule} and \cref{chain-rule},
\begin{equation}
	f'(x)=g'(x)\cos(h(x)) + g(x)(\cos'(h(x))\cdot h'(x)) \label{exercise1.1-fprime}
\end{equation}
First, $g'(x)$. Assuming $g(x)>0$ which is given:
\begin{align*}
	g(x)                                                                                       & = x^{\alpha}       \\[0.5em]
	\text{Take the natural log of both sides: }\ln(g(x))                                       & = \ln(x^{\alpha})  \\[0.5em]
	\ln(g(x))                                                                                  & = \alpha \ln(x)    \\[0.5em]
	\text{Differentiate both sides w.r.t. $x$: }(\ln(g(x)))'                                   & = (\alpha \ln(x))' \\[0.5em]
	\text{From \cref{chain-rule} and \cref{product-rule-for-constants}: }\ln'(g(x))\cdot g'(x) & = \alpha \ln'(x)
\end{align*}
\begin{equation}
	\text{So: }g'(x)=\dfrac{\alpha \ln'(x)}{\ln'(g(x))} \label{g'(x)}
\end{equation}
The derivative of the natural logarithm:
\begin{align*}
	\ln'(x) & = \lim_{h \to 0}\dfrac{\ln(x+h)-\ln(x)}{h} \text{ from the definition of the derivative.} \\[0.5em]
	        & = \lim_{h \to 0} \dfrac1h \ln\left( 1 + \dfrac{h}{x} \right)
\end{align*}
Let $ u= \dfrac{h}{x} \implies h = xu$
\begin{align*}
	\ln'(x) & = \lim_{u \to 0} \dfrac{1}{xu} \ln(1+u)             \\[0.5em]
	        & = \dfrac{1}{x} \lim_{u \to 0} \dfrac{1}{u} \ln(1+u) \\[0.5em]
	        & = \dfrac{1}{x} \lim_{u \to 0} \ln(1+u)^{\frac1u}    \\[0.5em]
	        & = \dfrac{1}{x} \ln(\lim_{u \to 0}(1+u)^{\frac1u})
\end{align*}
We do not need to evaluate this limit since it is, by definition, the base of the natural logarithm, $e$, and $\ln(e)=1$ from the definition
of the natural logarithm. So:
\begin{equation}
	\ln'(x)=\dfrac{1}{x} \label{ddx-of-ln}
\end{equation}
From \cref{g'(x)} and \cref{ddx-of-ln}:
\begin{equation}
	g'(x)=\dfrac{\dfrac{\alpha}{x}}{\dfrac{1}{g(x)}}=\dfrac{\alpha g(x)}{x}=\alpha x^{\alpha-1} \label{product-rule-for-any-power}
\end{equation}

Now, $h'(x)$:

Define $y(x) \coloneq x^{\beta}$ such that $h(x)$ = $\dfrac{1}{y(x)}$, then $h'(x)=\left(\dfrac{1}{y}\right)'(x)$ from \cref{ddx-of-1on-g}. Then:
\begin{equation}
	h'(x) = \dfrac{-y'(x)}{(y(x))^2}
\end{equation}
From \cref{product-rule-for-any-power}:
\begin{equation}
	h'(x)=\dfrac{-\beta x^{\beta-1}}{x^{2\beta}}=-\beta x^{-\beta-1} \label{exercise1.1-hprime}
\end{equation}

Since, the derivative of cosine is given as negative sine, \cref{exercise1.1-fprime} becomes, from \cref{product-rule-for-any-power} and
\cref{exercise1.1-hprime}
\begin{equation}
	\alpha x^{\alpha-1}\cos\left(\dfrac{1}{x^\beta}\right)+\beta x^{\alpha-\beta-1}\sin\left(\dfrac{1}{x^{\beta}}\right) \label{exercise1.1-result}
\end{equation}

\subsection{}
Determine all the values $\alpha,\beta \in (0,+\infty)$ such that
\begin{equation*}
	\lim_{x \to 0^+}\dfrac{f(x)-f(0)}{x} \text{ exists.}
\end{equation*}
\subsubsection{Solution}
From \cref{exercise1}:
\begin{equation}
	\lim_{x \to 0^+}\dfrac{f(x)-f(0)}{x} = \lim_{x \to 0^+}\dfrac{x^{\alpha}\cos(x^{-\beta})-0}{x}=\lim_{x \to 0^+}x^{\alpha-1}\cos\left(\dfrac{1}{x^{\beta}}\right)
\end{equation}

Case 1: $\alpha = 1 \implies \alpha-1=0 \implies x^{\alpha-1}=1$. The limit becomes $\lim_{x \to 0^+}\cos\left(\dfrac{1}{x^{\beta}}\right)$
which does not exist.

Proof:

Assume for contradiction that the limit does exist and equals $l$, then, $\forall \epsilon > 0, \exists \delta > 0$, such that if
$0< |x-a| < \delta, |f(x)-l|<\epsilon$. Choose a sufficiently large natural number $N$, such that $(\frac{1}{2\pi n})^{1/\beta}$ is less than $\delta$.
Then there will be a point $a=(\frac{1}{2\pi n})^{1/\beta}$ where $\cos(a)=1$ and a point $b=(\frac{1}{2\pi n+1})^{1/\beta}<a<\delta$ where
$\cos(b)=-1$.

Choose $\epsilon=1$. From the definition of the limit, $|f(a)-l|<1 \implies |1-l|<1$, and $|f(b)-l|<1 \implies | -1 -l|<1$.

From the triangle inequality:

$|1-(-1)| = |(1-l)+(l+1)| \le |1-l| + |1+l| < 1+1 < 2 \implies 2<2$. Which is a contradiction. Thus, the limit does not exist. \qed

Case 2: $0<\alpha < 1 \implies \alpha-1<0$ The term $x^{\alpha-1}$ goes to infinity as $x$ goes to 0, and since the $\cos(\frac{1}{x^\beta})$ term
oscillates between 0 and 1 and does not approach any limit at 0 as shown above, the limit diverges and does not exist.

Proof:

A function $f$ approaches $+\infty$ at $a$ if $\forall M \in \mathbb{R}, \exists \delta>0$ such that $\forall x$, if $0< |x-a| <\delta$, $f(x)>M$.
Let $p=\alpha - 1$. Since $0<\alpha<1, -1<p<0$. Let $q=-p$, then $0<q<1$. $x^p=x^{-q}=\frac{1}{x^q}$. If we choose $\delta=(\frac{1}{M})^{1/q}$, then
for all $0<x<\delta$, $x^p>M$.

To verify, choose $M$ to be arbitrary, and $\delta=(\frac{1}{M})^{1/q}=(\frac{1}{M})^{1/(1-\alpha)}$. Then $x<(\frac{1}{M})^{1/(1-\alpha)}$, so
$x^{1-\alpha}<\frac1M \implies x^{-(1-\alpha)} > M \implies x^{\alpha -1} > M$. \qed

Case 3: $\alpha > 1 \implies x^{\alpha -1} > 0$ when $x>0$. $x^{\alpha-1}$ goes to 0 as $x$ goes to 0 from the right. Since
$0 \le |\cos(\frac{1}{x^\beta})| \le 1, |x^{\alpha-1}\cos(\dfrac{1}{x^\beta})| \le x^{\alpha-1}$ which goes to 0 as $x \to 0^+$, so the whole
function goes to 0 as $x \to 0^+$ since the limit of products is the product of limits (\cref{limit-of-products}).

Proof:

From the definition of a limit, $\forall \epsilon>0, \exists \delta>0$, such that, if $0<x<\delta, |f(x)-0|<\epsilon$. $|f(x)|$ is bounded by $x^{\alpha-1}$ since
$0 \le |\cos(x)| \le 1, \forall x$. So if we choose $\epsilon > x^{\alpha-1}$, $x < \epsilon^{\frac{1}{\alpha-1}}$.
Thus, $\delta=\epsilon^{\frac{1}{\alpha-1}}$.

To verify, let $\epsilon > 0$ be given, and choose $\delta=\epsilon^{\frac{1}{\alpha-1}}$. For all $x$ such that $0<x<\delta$, $x<\epsilon^{\frac{1}{\alpha-1}}$.
Raise both sides to the power of $\alpha-1$, $x^{\alpha-1}<\epsilon$. Since $|x^{\alpha-1}\cos(\frac{1}{x^\beta})|\le x^{\alpha-1}<\epsilon$, it
follows that $|f(x)|<\epsilon$. \qed

Considering all the three cases together, the only values of $\alpha$ and $\beta \in (0,+\infty)$ where the limit exists are
$\alpha \in (1,+\infty), \beta \in (0,+\infty)$.

\subsection{}
Determine all values of $\alpha, \beta \in (0,+\infty)$ such that
\begin{equation*}
	\lim_{x \to 0^+}f'(x) \text{ exists.}
\end{equation*}

\subsubsection{Solution}
The derivative of $f(x), f'(x)$ (\cref{exercise1.1-result}) is
\begin{equation*}
	\underbrace{\alpha x^{\alpha-1}\cos\left(\dfrac{1}{x^\beta}\right)}_{\text{Term One}}+\underbrace{\beta x^{\alpha-\beta-1}\sin\left(\dfrac{1}{x^{\beta}}\right)}_{\text{Term Two}}
\end{equation*}

Each term behaves similarly and will oscillate between $-1$ and 1 as $x$ goes to 0 from the right. I have already proven above that for the limit
of term one to exist, $\alpha > 1$. I have also already proven that if $\alpha-\beta-1<0$, the limit will not exist since $x^{\alpha-\beta-1}$ will
diverge. So $\alpha-\beta-1$ must at least be zero. Since the sum of limits is the limit of sums (\cref{limit-of-sums}), we can consider each
term separately.

Case 1: $\alpha-\beta-1=0$ and the limit of term two becomes $\beta\sin\left(\dfrac{1}{x^\beta}\right)$. The
$\lim_{x \to 0^+}\beta\sin\left(\dfrac{1}{x^\beta}\right)$ does not exist for the same reason the limit of cosine of the same term does not exist.

Assume for contradiction that the limit does exist and equals $l$, then, $\forall \epsilon > 0, \exists \delta > 0$, such that if
$0< |x-a| < \delta, |f(x)-l|<\epsilon$. Choose a sufficiently large natural number $N$, such that $(\frac{1}{2\pi n})^{1/\beta}$ is less than $\delta$.
Then there will be a point $a=(\frac{1}{2\pi n+\pi/2})^{1/\beta}$ where $\beta\sin(a)=\beta$ and a point $b=(\frac{1}{2\pi n+3\pi/2})^{1/\beta}<a<\delta$ where
$\beta\sin(b)=-\beta$.

Choose $\epsilon=\beta$. From the definition of the limit, $|f(a)-l|<\beta \implies |\beta-l|<\beta$, and $|f(b)-l|<\beta \implies | -\beta -l|<\beta$.

From the triangle inequality:

$|\beta-(-\beta)| = |(\beta-l)+(l+\beta)| \le |\beta-l| + |\beta+l| < \beta+\beta < 2\beta \implies 2<2 \text{ (since } \beta>0)$.
Which is a contradiction. Thus, the limit does not exist. \qed

Case 2: $\alpha-\beta-1>0$. $x^{\alpha-\beta-1}$ goes to 0 as $x$ goes to 0 from the right. Since $0 \le |\sin(\frac{1}{x^\beta})| \le 1$,
$|\beta x^{\alpha-\beta-1}\sin(\frac{1}{x^\beta})| \le \beta x^{\alpha-\beta-1}$ which goes to 0 as $x \to 0^+$, so the whole function goes to 0
as $x \to 0^+$ since the limit of products is the product of limits (\cref{limit-of-products}).

Proof:

From the definition of a limit, $\forall \epsilon > 0, \exists \delta > 0$, such that, if $0 < x < \delta, |f(x)-0| < \epsilon$. $|f(x)|$ is bounded
by $\beta x^{\alpha-\beta-1}$ since $0 \le \sin(x) \le 1, \forall x$. So if we choose $\epsilon > \beta x^{\alpha-\beta-1}$,
$x < (\frac{\epsilon}{\beta})^{\frac{1}{\alpha-\beta-1}}$. Thus, $\delta = (\frac{\epsilon}{\beta})^{\frac{1}{\alpha-\beta-1}}$.

To verify, let $\epsilon > 0$ be given, and choose $\delta = (\frac{\epsilon}{\beta})^{\frac{1}{\alpha-\beta-1}}$. For all $x$ such that, if
$0 < x < \delta$, $x < (\frac{\epsilon}{\beta})^{\frac{1}{\alpha-\beta-1}}$. Raise both sides to the power of $\alpha-\beta-1$, so
$x^{\alpha-\beta-1}<\frac{\epsilon}{\beta} \implies \beta x^{\alpha-\beta-1}<\epsilon$.
Since $|\beta x^{\alpha-\beta-1}\sin(\frac{1}{x^\beta})| \le \beta x^{\alpha-\beta-1}<\epsilon$, this implies $|f(x)|<\epsilon$.

Considering the solution to 1.2, as well as the cases where the limit of the second term exists, we require that $\alpha-\beta-1>0$ and $\alpha>1$.
So, $\alpha > \beta + 1$ which requires $\beta > 0$ to satisfy $\alpha>1$ for the limit of the derivative to exist as $x$ goes to 0 from the right.

\section{Exercise 2}
We say that $f : \mathbb{R} \to \mathbb{R}$ is even if for every $x \in \mathbb{R}$ it holds that $f(x)=f(-x)$. Also, we say that $f$ is odd if
for every $x \in \mathbb{R}$ it holds that $f(x)=-f(-x)$. Prove that if $f$ is differentiable everywhere in $\mathbb{R}$ then
\begin{equation*}
	f' \text{ is even if }f \text{ is odd, and that }f'\text{ is odd if }f \text{ is even.}
\end{equation*}
Assume that $f$ is $k$-times differentiable. What can be said about $f^{(k)}$?

\subsection{Solution}

Suppose that $f$ is even, so $f(x)=f(-x)$. Differentiating both sides yields $f'(x)=f'(-x)\cdot(-x)'$ from \cref{chain-rule}. The derivative of $-x$
\begin{align*}
	(-x)' & = \lim_{h \to 0}\dfrac{-(x+h)-(-x)}{h} \\[0.5em]
	      & = \lim_{h \to 0}\dfrac{-x-h+x}{h}      \\[0.5em]
	      & = \lim_{h \to 0}\dfrac{-h}{h}          \\[0.5em]
	      & = -1
\end{align*}
So $f'(x)=-f'(-x)$ which is odd by definition. So the first derivatives of even functions are odd. \qed

Suppose that $g$ is odd, so $g(x)=-g(-x)$. Differentiating both sides yields $g'(x)=(-1)g'(-x)(-x)'$ from \cref{chain-rule} and \cref{product-rule-for-constants}.
The derivative of $-x$ has already been shown to be $-1$, so $g'(x)=(-1)(-1)g'(-x)=g'(-x)$ which is even by definition. So the first derivative of
odd functions are even. \qed

It is a direct corollary of the above two statements that, if the first derivative of an even function is odd, and the first derivative of an odd function, even, then
the second derivative of an even function is even, and the second derivative of an odd function, odd. Thus, for the $k'th$ derivative of $f$, if $k$ is
even, then $f^{(k)}$ will be the same as $f$ in terms of being even or odd, and if $k$ is odd, then $f^{(k)}$ will be opposite $f$ in terms of being
even or odd.

\section{Exercise 3}
Suppose that $f$ is a continuous function on the reals and
\begin{equation*}
	\lim_{x \to +\infty}\dfrac{f(x)}{x^{27}}=\lim_{x \to -\infty}\dfrac{f(x)}{x^{27}}=0
\end{equation*}
Prove that there exists $x_0$ such that $x_0^{27}+f(x_0)=0$

\subsection{Solution}

Define $g(x) \coloneq x^{27}+f(x)$.

Since the $\lim_{x \to +\infty}\dfrac{f(x)}{x^{27}}=0$, there exists some $M \in \mathbb{R^+}$ such that
for all $x > M$, $\left| \dfrac{f(x)}{x^{27}}\right|<\dfrac{1}{2}$ from the definition of a limit (choosing $\epsilon=1/2$). It follows that
$|f(x)| < \dfrac{x^{27}}{2}$. Take some $a>M>0$, then $g(a)=a^{27}+f(a) > a^{27}-\dfrac{a^{27}}{2}=\dfrac{a^{27}}{2}>0$.

Since the $\lim_{x \to -\infty}\dfrac{f(x)}{x^{27}}=0$, there exists some $M' \in \mathbb{R^+}$ such that for all
$x < -M'$, $\left| \dfrac{f(x)}{x^{27}}\right|<\dfrac{1}{2} \implies |f(x)|<\dfrac{|x|^{27}}{2}$.
Take some $b < -M' < 0$, then $g(b)=b^{27}+f(b) < b^{27}+\dfrac{|b|^{27}}{2}=b^{27}-\dfrac{b^{27}}{2}<0$. Note that $|b|^{27}=-b^{27}$ since $27$ is
odd.

Since $x^{27}$ is continuous and $f$ is continuous, $g$ is also continuous on the reals (\cref{continuous-sum-is-continuous})
and so is continuous on $[b,a]$. Since $f(b)<0<f(a)$ then by
\cref{mean-value} there is some $x_0 \in [b,a]$ where $g(x_0)=0$ and therefore some point where $x_0^{27}+f(x_0)=0$. \qed

\section{Exercise 4}
Let $f$ and $g$ be bounded functions on the reals. Suppose that $f$ and $g$ are uniformly continuous and for all $x \in \mathbb{R}$ let
\begin{equation}
	h(x) \coloneq f(x)g(x)
\end{equation}
Prove that $h$ is uniformly continuous. If we drop the assumption that $f$ and $g$ are bounded, is $h$ necessarily uniformly continuous?

\subsection{Solution}
It is given that $f$ and $g$ are bounded, so there must be some $N,M \in \mathbb{R}$ such that $|f(x)|\le N$, and $|g(x)|\le M$. To show $h$ is
uniformly continuous, we must show that $\forall \epsilon > 0, \exists \delta >0$ such that $\forall x,y \in \mathbb{R}$, if $|x-y|<\delta, |h(x)-h(y)|<\epsilon$.
\begin{align*}
	|h(x)-h(y)| & = |f(x)g(x)-f(y)g(y)|                         \\[0.5em]
	            & = |f(x)g(x) - f(y)g(x) + f(y)g(x) - f(y)g(y)| \\[0.5em]
	            & \le |f(x)g(x)-f(y)g(x)| + |f(y)g(x)-f(y)g(y)| \\[0.5em]
	            & = |g(x)||f(x)-f(y)| + |f(y)||g(x)-g(y)|       \\[0.5em]
	            & \le N|f(x)-f(y)|+M|g(x)-g(y)|
\end{align*}
Let $\epsilon>0$. Since $f$ is uniformly continuous, $\exists \delta_1 >0$ such that
\begin{equation*}
	|x-y| < \delta_1 \implies |f(x)-f(y)|<\dfrac{\epsilon}{2N}
\end{equation*}
Since $g$ is uniformly continuous, $\exists \delta_2 >0$ such that
\begin{equation*}
	|x-y| < \delta_2 \implies |g(x)-g(y)|<\dfrac{\epsilon}{2M}
\end{equation*}
If we choose $\delta = \text{min}(\delta_1,\delta_2)$ then when
\begin{equation*}
	|x-y| < \delta \implies |h(x)-h(y)| \le N|f(x)-f(y)|+M|g(x)-g(y)| < N \cdot \dfrac{\epsilon}{2N} + M \cdot \dfrac{\epsilon}{2M}=\epsilon
\end{equation*}
Which is the definition of uniformly continuous, therefore $h(x)$ is uniformly continuous. \qed

If we drop the assumption that $f$ and $g$ are bounded, then their product is not necessarily uniformly continuous.

Consider $f(x)=g(x)=x$ which is uniformly continuous but not bounded.

Proof:

Choose $\delta = \epsilon$, then when $|x-y|<\delta, |f(x)-f(y)| = |x-y| < \delta = \epsilon$.

Their product, $h(x)=x^2$ is not uniformly continuous. Choose $\epsilon=1$ and let $\delta>0$.

Consider $x=\frac{1}{\delta}$ and $y=\frac{1}{\delta}+\frac{\delta}{2}$. $|x-y| = |\frac{1}{\delta}-(\frac{1}{\delta}+\frac{\delta}{2})| = \frac{\delta}{2}<\delta$.

$|f(x)-f(y)| = |x^2-y^2| = |x-y||x+y| = |\frac{\delta}{2}||\frac{1}{\delta}+(\frac{1}{\delta}+\frac{\delta}{2})| = (\frac{\delta}{2})(\frac{2}{\delta}+\frac{\delta}{2})$
$=1+\frac{\delta^2}{4}>\epsilon$.

So $x^2$ is not uniformly continuous since $\exists \epsilon>0$ such that $\exists \delta > 0$ with $|x-y|<\delta$
but $|f(x)-f(y)>\epsilon$. \qed

\section{Exercise 5}
Find $f'(0)$ if
\begin{equation}
	f(x) \coloneq
	\begin{cases}
		\dfrac{g(x)}{x} & x \ne 0 \\
		0               & x=0
	\end{cases}
\end{equation}
and $g(0)=g'(0)=0$ and $g''(0)=17$.

\subsection{Solution}

From the definition of the derivative
\begin{align*}
	f'(0) & = \lim_{h \to 0}\dfrac{f(x+h)-f(x)}{h}                            \\[0.5em]
	      & = \lim_{h \to 0}\dfrac{f(0+h)-f(0)}{h}                            \\[0.5em]
	      & = \lim_{h \to 0}\dfrac{(g(h)/h)-0}{h}                             \\[0.5em]
	      & = \lim_{h \to 0}\dfrac{g(h)}{h^2} \text{ which is indeterminate.} \\[0.5em]
\end{align*}
From \cref{le-medical-centre} and \cref{power-rule}:
\begin{equation}
	\lim_{h \to 0}\dfrac{g(h)}{h^2}=\lim_{h \to 0}\dfrac{g'(h)}{2h}=\lim_{x \to 0}\dfrac{g''(h)}{2}=\dfrac{17}{2}
\end{equation}

\end{document}
